\chapter*{Abstract}

This work describes the development and implementation of a comprehensive calibration
and optimization protocol for an open-architecture X-ray computed microtomography
($\mu$CT) system. The central objective was to transform a set of independent components
into a high-precision instrument for scientific research, comparable to a commercial
device.

The methodology began with hardware stabilization, securing the source and the detector
onto a rigid base to eliminate accidental misalignments. Physical alignment techniques
were implemented through the analysis of cone-beam projections and the superimposition
of cylinder projections, successfully reducing angular deviations to less than 5$^{\circ}$
and bringing the principal point of the beam closer to the center of the detector.
During the analytical calibration stage, two methods were evaluated: a multiple-projection
method, analyzing the elliptical trajectories of metallic spheres, which proved to
be the most robust and consistent for determining fine geometric parameters; and a
single-projection method, based on the relationship between the sides of a projected
quadrilateral, which was discarded due to its extreme sensitivity to positioning
errors.

Furthermore, the effective spatial resolution of the system was characterized based
on a modulation transfer function (MTF) study. The results placed the resolution
between 15 and 60 $\mu$m, and the empirical curve as a function of magnification
was obtained. Finally, digital pre-processing techniques such as flat-field correction
(FFC) and temporal uniformity correction were applied, which eliminated electronic
patterns from the detector, increased the signal-to-noise ratio (SNR), and standardized
the beam intensity in area and time for a set of projections. The protocol was
successfully validated through the three-dimensional reconstruction of biological
and mechanical samples with high sharpness and the absence of artifacts, achieving
a microtomography system with results equivalent to those of a commercial unit.